% ============================================================
%  Chapter 4 – System Architecture
% ============================================================

\section{System Architecture}

\subsection{Design Principles and Requirements}
\label{sec:design-principles}

The architecture of the system was conceived around a set of normative requirements that
distinguish it from general-purpose question-answering systems. These requirements reflect
both the domain specificity of parliamentary discourse and the democratic accountability
imperatives that motivate the research.

\paragraph{Balanced political representation.}
Parliamentary debate is inherently multi-stakeholder. A system that privileges the majority
coalition or the most vocally active deputies would reproduce the very asymmetries that
limit journalistic and civic access to legislative information. The architecture therefore
treats proportional coverage across parliamentary groups as a first-class design
constraint, embedded not only in the retrieval stage but also in the generation and
evaluation pipelines. Coverage is measured quantitatively and enforced algorithmically at
multiple stages of the processing pipeline, preventing any single ideological faction from
dominating the synthesised response.

\paragraph{Multi-view answer generation.}
Rather than producing a single, consensus-seeking response, the system is designed to make
political divergence explicit. This principle, which may be termed \textit{deliberative
transparency}, requires the generation module to preserve and articulate positions that are
in tension, attributing claims to their sources and refraining from premature synthesis.
The architectural implication is that generation cannot be a monolithic operation: it
requires a decomposition into party-specific narrative units before any integrative step is
attempted.

\paragraph{Citation faithfulness.}
The evidentiary chain between a generated claim and its source in parliamentary record must
be verifiable and exact. The system adopts an \textit{offset-based citation model}, in
which each quotation is extracted from its source text using byte-level character positions
stored at ingestion time. This design choice makes it structurally impossible for the
system to generate citations that are paraphrases or semantic approximations: the quoted
string is always a contiguous substring of the original parliamentary speech. Faithfulness
is not treated as a post-hoc evaluation criterion but as an architectural invariant.

\paragraph{Query-dependent authority modeling.}
The notion of expert authority cannot be defined context-independently. A deputy with
extensive legislative activity on labour law is not necessarily authoritative when the
query concerns environmental regulation. The system therefore computes authority scores
as a function of both the speaker's profile and the semantic content of the incoming query.
This requires that authority computations be deferred to query time and that they operate
on a shared embedding space that couples query meaning with speaker expertise signals.

\paragraph{Interpretability.}
Downstream users of the system---journalists, researchers, parliamentary monitors---require
not only an answer but an understanding of how that answer was formed. The architecture
accordingly exposes the ideological positioning of speakers, the authority basis for
citation choices, and the provenance of every quoted passage. Interpretability is thus a
structural property of the output format, not an optional add-on.

\paragraph{Modularity.}
Each functional concern---retrieval, authority scoring, ideological projection, generation,
citation validation---is encapsulated in a distinct component with well-defined interfaces.
This modularity serves multiple purposes: it enables independent evaluation of each stage,
facilitates targeted improvement without global refactoring, and allows the system to be
extended with new components without disrupting existing ones.


\subsection{High-Level Overview}
\label{sec:high-level-overview}

The system is organised as a \textit{layered architecture} comprising three principal
tiers---presentation, application, and data---supplemented by a set of external services
that provide semantic intelligence. This stratification reflects a deliberate separation
between concerns of interaction, processing, and persistence, and it ensures that changes
in any one tier do not propagate unpredictably into the others.

At the highest level, a user submits a natural language query about parliamentary activity.
The query traverses the application layer, where it triggers a cascade of retrieval,
scoring, ideological analysis, and generation operations. The result is a structured
multi-view response, grounded in cited parliamentary speeches, which is then rendered in
the presentation layer together with interpretive visualisations of the underlying evidence.

\subsubsection{Presentation Layer}

The presentation layer is responsible for mediating between the complexity of the
underlying analytical pipeline and the cognitive needs of the end user. Its primary
function is not merely to display a textual answer but to render the epistemic structure of
that answer---to make visible the sources of claims, the identities and affiliations of
cited speakers, their positions in the ideological space, and the degree of consensus or
divergence among parliamentary groups.

Multi-view answers are presented with per-party attribution, allowing the user to navigate
political positions in a structured manner rather than consuming a homogenised synthesis.
Each cited passage is displayed alongside metadata identifying the speaker, their
parliamentary group, the session date, and a unique identifier that enables traceability
to the original parliamentary record. An ideological compass visualisation provides a
data-driven, two-dimensional representation of the positioning of the evidence fragments
relative to the query, offering a geometric perspective on the political landscape
underlying the generated answer.

The presentation layer additionally exposes an evaluation interface that allows operators
to assess system performance on a predefined set of queries, comparing the system's
responses against a curated baseline using both automated metrics and human preference
surveys. This evaluation component is architecturally integrated with the main pipeline
rather than isolated as a separate tool.

\subsubsection{Application Layer}

The application layer constitutes the computational core of the system and is itself
organised into five cooperating modules.

The \textit{retrieval engine} is responsible for selecting the evidence fragments from the
parliamentary corpus that are most relevant to the incoming query. It operates across two
complementary channels---one based on dense vector similarity, the other on structured
graph traversal---and combines their outputs through a multi-factor scoring function that
jointly considers semantic relevance, speaker authority, political diversity, and opinion
salience. The engine includes a query rewriting sub-module that expands short queries
before embedding, and a neighbourhood expansion mechanism that enlarges retrieved fragments
to preserve discourse context.

The \textit{authority scoring module} computes, for each candidate speaker in the retrieved
evidence, a composite authority score that reflects the speaker's expertise relative to the
query's semantic domain. The score is the weighted sum of six components corresponding to
professional background, educational history, committee membership, legislative
co-signatorship, speech-level topical engagement, and institutional role. All components
are normalised to a common scale and computed using the query embedding as a semantic
reference, making authority a dynamic, context-sensitive quantity.

The \textit{ideological modeling module} projects the retrieved evidence fragments onto a
two-dimensional space derived from the principal axes of variance in the set's semantic
embeddings. The resulting compass representation captures ideological proximity and
divergence among parliamentary groups in a data-driven, query-specific manner, without
relying on fixed party labels as the primary analytic unit.

The \textit{multi-stage generation module} orchestrates a sequential pipeline of
large language model calls that progressively transform the retrieved evidence into a
structured multi-view answer. Each stage has a distinct functional role: decomposing the
query into atomic claims, generating party-specific narrative sections, integrating those
sections into a coherent discourse, and enforcing verbatim citation insertion. The
sequential design reflects the observation that a single-pass generation cannot
simultaneously satisfy the competing constraints of factual attribution, political
balance, and narrative coherence.

The \textit{citation validation module} operates as a cross-cutting concern that enforces
the faithfulness invariant across all stages of the pipeline. It maintains a registry of
citation expectations, audits the output of each generation stage, and triggers repair
mechanisms when citations are lost or corrupted during integration.

\subsubsection{Data Layer}

The data layer is built around a \textit{property graph database} that models the entities
and relationships of the parliamentary domain. The graph structure is not incidental: it
enables the efficient traversal of relationships---between deputies, their committees, the
acts they co-signed, the sessions in which they spoke, and the parties they represent at
any given point in time---that are essential to both retrieval and authority scoring.

Parliamentary speeches are stored in their full textual form, with a separate collection of
\textit{chunks} derived from each speech by a fixed-length sliding-window procedure. Each
chunk carries a dense vector embedding in a high-dimensional semantic space, enabling
approximate nearest-neighbour retrieval. Crucially, each chunk also stores the exact
byte-level character offsets of its boundaries within the original speech text, forming
the foundation of the citation faithfulness mechanism.

Speaker metadata is stored as node attributes, including pre-computed embeddings of
professional and educational descriptions, which allow authority components to be computed
efficiently at query time without re-embedding static profile information. Parliamentary
acts are stored with their thematic descriptions and pre-computed embeddings, supporting
the semantic evaluation of a deputy's legislative engagement relative to a given query
topic.

The graph is supplemented by vector indices that allow semantic retrieval to be performed
natively within the same infrastructure, avoiding the need for a separate vector store and
preserving the ability to combine structured graph traversal with unstructured semantic
search in a single query operation.

\subsubsection{External Services}

The system integrates two categories of external services: language model providers and
institutional data sources.

Large language models are invoked exclusively within the generation module and serve as the
primary mechanism for natural language synthesis. The system uses a state-of-the-art
generative model---specifically GPT-4o---across all generation stages, with a uniformly
low temperature setting to favour factual fidelity over creative elaboration. The choice
of a single model across all stages simplifies dependency management and ensures
consistency in output style.

The embedding service---text-embedding-3-small, a 1536-dimensional dense encoder---is the
shared representation backbone of the entire system. Every semantic similarity computation
in the pipeline, from retrieval scoring to authority component evaluation to ideological
projection, operates on embeddings produced by this model. Its consistent use across all
modalities ensures that the semantic space is coherent and that comparisons between query
embeddings, chunk embeddings, and speaker profile embeddings are geometrically meaningful.

Institutional data sources supply the raw parliamentary record and the structured metadata
that the system ingests. These sources are discussed in detail in the following subsection.


\subsection{Data Acquisition and Preprocessing}
\label{sec:data-acquisition}

\subsubsection{Institutional Data Sources}

The primary data source is the official open data infrastructure of the Italian Parliament,
which publishes machine-readable records of parliamentary sessions, including full
transcripts of plenary and committee debates, structured information on legislative acts
and their signatories, deputy biographical data, and committee membership records with
associated temporal validity intervals. These records are provided under an open licence
and updated incrementally as parliamentary activity progresses.

The transcripts represent the evidential core of the system. Each transcript is associated
with a session date, a debate identifier, and a phase marker that captures the structural
organisation of parliamentary proceedings. Deputy profiles contain professional and
educational history in free-text form, which must be embedded before they can be used
in authority computations. Legislative acts carry thematic descriptors that are similarly
embedded to support semantic matching.

A secondary thematic resource is a curated ontology of committee topics, which maps each
parliamentary committee to a set of domain-specific keywords representing its areas of
competence. This mapping supplements the structural committee membership data with
semantic content, enabling authority scoring to assess the topical relevance of a
deputy's committee roles with greater precision than committee name alone would permit.

\subsubsection{Text Preprocessing and Alignment Map}

The raw parliamentary transcripts require substantial preprocessing before they are suitable
for indexing. The primary challenge is that transcripts contain a mixture of substantive
speech content and procedural markers---presiding officers' announcements, voting
declarations, quorum calls, formal greetings, and other parliamentary formalities---that
carry no informational value relative to policy queries. These procedural elements are
identified and removed through a combination of structural heuristics and textual pattern
matching, leaving a cleaned representation that preserves the substantive content of each
deputy's interventions.

Text normalisation is applied to resolve encoding inconsistencies, Unicode variation forms,
and whitespace irregularities that arise from the heterogeneity of source document
formats. The normalisation procedure produces a canonical representation of each speech
that is consistent across the corpus.

A critical design decision concerns the relationship between the cleaned representation used
for indexing and retrieval, and the original raw text used for citation extraction. Rather
than discarding the raw text after preprocessing, the system maintains an
\textit{alignment map} for each speech that records, for every position in the cleaned
text, the corresponding byte offset in the original source. This bidirectional mapping
allows the system to retrieve evidence fragments using the processed text while extracting
citations directly from the verbatim original. The alignment map is the architectural
mechanism that makes citation faithfulness a structural guarantee: no post-hoc
string-matching or semantic verification of quoted passages is necessary, because the
quoted string is retrieved deterministically from a stored offset pair.

\subsubsection{Chunking Strategy}

Individual parliamentary speeches can span many thousands of characters, rendering the full
text unsuitable as a retrieval unit for two reasons. First, a single embedding of the
entire speech would suppress the topical heterogeneity that characterises lengthy
parliamentary interventions, in which a deputy may address several distinct policy areas in
sequence. Second, the granularity of citations would be insufficient: a passage-level
quotation is more informative and verifiable than a speech-level reference.

The system therefore adopts a \textit{sentence-aware sliding window} approach to
segmentation. Speeches are divided into overlapping chunks of approximately fixed character
length, with boundaries adjusted to coincide with sentence boundaries so that no chunk
begins or ends mid-sentence. The overlap between consecutive chunks---set to a fraction of
the maximum chunk length---ensures that passages occurring near chunk boundaries are
represented in at least two retrieval units, mitigating the risk that relevant content is
artificially split across non-retrievable gaps.

Each chunk is assigned a dense embedding computed over its textual content, and its
start and end offsets within the original speech are stored explicitly. A directed
\textit{next-chunk} relationship connects consecutive chunks within the same speech,
enabling the retrieval engine to expand a retrieved chunk by appending its immediate
successor when discourse context requires it.

The choice of chunk length involves a trade-off between retrieval precision and context
preservation. Shorter chunks are more likely to contain a single coherent statement, but
they may lack the context needed for an LLM to interpret the statement correctly.
Longer chunks preserve more context but dilute the embedding signal, reducing retrieval
precision. The adopted parameters represent an empirically motivated balance between these
competing considerations, validated through retrieval quality experiments on the
evaluation set.

\subsubsection{Build and Update Pipeline}

The construction of the knowledge graph and its associated vector indices is decomposed
into a \textit{build phase} and an \textit{incremental update phase}, each with distinct
operational characteristics.

The full build phase ingests the complete parliamentary corpus from scratch, constructing
all graph nodes and relationships, computing and storing all embeddings, and creating the
vector indices. This phase is computationally intensive---the embedding computation alone
requires a large number of API calls to the embedding service---and is therefore executed
offline, prior to any runtime query processing. Authority component embeddings for
speaker profiles and legislative act descriptions are precomputed and stored during this
phase, ensuring that query-time authority scoring can proceed without incurring additional
embedding latency.

The incremental update phase processes new sessions and acts as they become available from
the institutional data sources. Because parliamentary records are published with a delay
following each session, the update frequency is bounded by the publication schedule of the
source. Incremental updates add new nodes and relationships without modifying existing
ones, preserving the consistency of the graph while extending its temporal coverage.

An important architectural decision is the strict separation between the ingestion and
runtime components. The ingestion pipeline has write access to the graph database and is
responsible for all data transformations; the runtime query pipeline has only read access.
This separation prevents accidental data corruption during query processing and allows the
ingestion pipeline to be run independently---or even by a different operational role---
without affecting system availability.


\subsection{Data Model and Knowledge Graph Schema}
\label{sec:data-model}

\subsubsection{Node Types}

The knowledge graph is structured around a set of entity types that correspond to the
principal actors and artefacts of parliamentary life.

\textit{Deputy} nodes represent elected members of parliament and carry biographical
attributes---name, profession, educational history---alongside pre-computed embeddings of
the free-text professional and educational fields. These embeddings are the primary
mechanism through which expertise signals are integrated into authority scoring.

\textit{GovernmentMember} nodes share the structural role of the Deputy type but represent
members of the executive branch---ministers, vice-ministers, and undersecretaries---who
intervene in parliamentary proceedings in an institutional rather than representative
capacity. Their authority is modeled through a role-based component that reflects the
institutional weight of their position.

\textit{ParliamentaryGroup} nodes encode the political formations to which deputies belong,
and their composition changes over time as deputies switch allegiances. Because group
membership is inherently temporal, the relationship between a deputy and their group
carries explicit start and end date attributes, enabling the system to resolve party
affiliation at any historical reference point.

\textit{Speech} nodes represent individual parliamentary interventions, each linked to a
specific session phase and attributed to a single speaker. The speech carries the full
preprocessed text of the intervention and is the parent entity from which chunk nodes are
derived.

\textit{Chunk} nodes are the primary retrieval unit. Each chunk stores its textual content
in the form used for display, its dense vector embedding, and the byte-level offsets that
map its boundaries back to the original speech text. The chunk is thus simultaneously an
index unit, a preview representation, and a citation anchor.

\textit{ParliamentaryAct} nodes encode legislative proposals, amendments, and motions.
Each act carries a thematic description and pre-computed semantic embeddings that allow
its topical domain to be compared against query embeddings during authority scoring.
Signatorship relationships connect acts to the deputies who co-sponsored them, with a
role attribute distinguishing primary from co-sponsors.

\textit{Session}, \textit{Debate}, and \textit{Phase} nodes form a hierarchical schema
that captures the structural organisation of parliamentary proceedings. A session
encompasses one or more debates; each debate is subdivided into phases; each phase
contains the speeches delivered within that procedural context. This hierarchy enables
structured retrieval queries that filter evidence by temporal or procedural criteria.

\textit{Committee} nodes represent parliamentary committees and serve as the bridge between
deputy membership records and the thematic keyword ontology. Their authority relevance is
determined not only by name but by the domain-specific keywords associated with each
committee, which are embedded and compared against query embeddings to produce fine-grained
topical relevance scores.

\subsubsection{Relationship Types}

The graph schema captures several distinct categories of relationship.

Biographical relationships connect deputies to their committees, parties, and institutional
roles, all with temporal validity attributes. The temporality of these relationships is
essential: a deputy who chaired a finance committee in a previous legislature cannot be
credited with that role when answering queries in the current legislative context unless
the query explicitly concerns historical events. Without a graph model that supports
temporal relationship attributes, this kind of validity-aware reasoning would require
either materialising multiple copies of the graph for different time slices or performing
expensive post-hoc filtering.

Discursive relationships connect speakers to their interventions and interventions to the
parliamentary session structure. These relationships enable the retrieval engine to
traverse from a retrieved chunk to its parent speech, from the speech to its speaker, and
from the speaker to their group affiliation---all in a single structured query.

Legislative relationships connect deputies to the parliamentary acts they have sponsored.
The distinction between primary and co-signatorship is preserved in the relationship type,
reflecting the fact that primary sponsors have a stronger claim to topical authority than
co-signatories.

Structural relationships among chunk nodes---specifically the directed \textit{next-chunk}
edge---model discourse continuity within a speech, enabling the neighbourhood expansion
logic described in the processing pipeline.

The graph formalism is appropriate for this domain because the informational value of the
data is encoded not only in entity attributes but in the structure of connections among
entities. A relational schema could represent individual entities efficiently but would
require expensive multi-table joins to answer the traversal queries that are central to
both retrieval and authority scoring. The graph model makes these traversals native and
efficient, while the integration of vector indices within the same infrastructure allows
semantic and structural retrieval to be combined without data movement between systems.

\subsubsection{Indexing Strategy}

Three categories of index support the retrieval operations of the runtime pipeline.

A \textit{vector index} is maintained over the chunk embedding attribute, enabling
approximate nearest-neighbour search in the dense embedding space. This index is the
primary mechanism for semantic retrieval: given a query embedding, the index returns the
set of chunk nodes whose embeddings have the highest cosine similarity to the query. The
vector index is constructed and updated as part of the ingestion pipeline and is queried
directly by the dense retrieval channel at runtime.

\textit{Property indexes} support exact and range lookups on structured attributes such as
session date, speaker identifier, and group name. These indexes are invoked by the graph
traversal channel and by the temporal validity filtering operations in authority scoring.
They also underpin the efficient resolution of cited speaker identities at generation time.

A \textit{full-text index} over thematic descriptions and act titles enables keyword-based
matching in the graph channel, complementing the semantic search with lexical precision
for proper nouns and domain-specific terms that may not be well-represented in the
embedding space.

The combination of these three index types---vector, property, and full-text---within a
single graph infrastructure supports \textit{hybrid retrieval}, in which semantic and
structural constraints are applied jointly. This hybrid strategy is central to the system's
ability to retrieve evidence that is both semantically relevant and structurally grounded
in the parliamentary record.


\subsection{Processing Pipeline}
\label{sec:processing-pipeline}

The runtime processing pipeline transforms a natural language query into a structured,
multi-view, citation-grounded answer through a sequence of nine interdependent stages.
Each stage has a well-defined input-output contract, and the pipeline is designed so that
the invariants established in early stages---in particular, the set of retrieved evidence
fragments and their authority scores---are preserved and propagated faithfully through the
later stages.

\paragraph{Stage 1: Query Embedding.}
The incoming query is first subjected to a rewriting step that expands underspecified or
short queries into a more descriptive form, while leaving sufficiently detailed queries
unchanged. The rewritten query is then encoded into a dense vector representation using
the shared embedding model. This query embedding serves as the universal semantic
reference for all subsequent stages: it is used in vector similarity search, in every
component of authority scoring, in ideological projection, and in the citation validation
step. The consistency of the embedding space across all these uses is what makes it
possible to define authority as a function of query semantics rather than as a fixed
property of the speaker.

\paragraph{Stage 2: Dual-Channel Retrieval.}
Evidence is collected through two complementary retrieval channels that are executed in
parallel and subsequently merged.

The \textit{dense channel} queries the vector index with the query embedding to retrieve
a large candidate pool of chunks---typically on the order of several hundred---with cosine
similarity above a minimum threshold. This channel is effective at capturing thematically
relevant passages regardless of the specific lexical forms used in either the query or
the source text.

The \textit{graph channel} performs a structured traversal of the knowledge graph,
identifying legislative acts whose thematic embeddings are semantically similar to the
query, and then traversing from those acts to the speeches delivered by their signatories.
This channel introduces a structural prior that favours speakers who have demonstrated
legislative commitment to the query's domain---not merely those who happen to have used
topically similar language in debate.

The two channels produce overlapping but complementary sets of evidence fragments, and
their union is substantially more diverse than either channel alone.

\paragraph{Stage 3: Re-ranking and Merging.}
The combined candidate pool is scored by a multi-factor function that simultaneously
considers five dimensions: semantic relevance (the cosine similarity of the chunk embedding
to the query), speaker authority (computed in Stage 4 and applied retroactively in the
final merge), political diversity (a penalty for over-representing any single
parliamentary group), opinion salience (an estimate of the informational density and
attitudinal specificity of the passage), and coverage balance (a correction term that
rewards fragments from underrepresented groups).

A \textit{salience expansion} mechanism detects chunks that are rhetorically weak---for
instance, brief procedural acknowledgements near the beginning of a speech---and expands
them by appending the immediately following chunk, which typically contains the substantive
content of the intervention. This expansion uses the \textit{next-chunk} graph
relationship and is applied before final scoring.

If the merged pool still under-represents certain parliamentary groups after initial
scoring, a \textit{coverage fill} operation executes targeted vector queries constrained
to the underrepresented groups, ensuring that the final evidence set provides adequate
representation across the ideological spectrum before generation begins.

\paragraph{Stage 4: Authority Scoring.}
For each speaker appearing in the retained evidence set, the authority scoring module
computes a composite score reflecting the speaker's domain expertise relative to the
query. The score integrates six components, each normalised to the unit interval.

The \textit{profession} and \textit{education} components assess the semantic proximity
between the speaker's professional and educational background---encoded as pre-stored
embeddings---and the query embedding. A speaker whose declared expertise is closely
aligned with the query domain receives a higher score on these components.

The \textit{committee} component evaluates the topical relevance of the speaker's current
and past committee memberships, using a keyword-enriched thematic ontology rather than
raw committee names. Only memberships that are temporally valid at the query reference
date are considered. The component takes the maximum relevance across all active
memberships, reflecting the principle that a speaker's most relevant committee defines
their authority on that topic.

The \textit{acts} component aggregates the speaker's legislative activity by summing the
relevance-weighted, time-discounted contributions of each parliamentary act they have
co-signed, with primary signatorship weighted more heavily than co-signatorship. Acts
whose thematic embeddings are insufficiently similar to the query are excluded entirely,
ensuring that prolific but off-topic legislative activity does not inflate the score.

The \textit{interventions} component applies the same aggregation logic to the speaker's
speech history, using the embeddings of individual speeches as the relevance signal and a
shorter time-decay half-life that reflects the greater temporal sensitivity of verbal
interventions relative to legislative commitments.

The \textit{role} component assigns a fixed base weight to speakers occupying formal
institutional roles---committee presidents, vice-presidents, secretaries, and government
members---with an optional bonus when the institutional role is semantically aligned with
the query. This component ensures that institutional authority is reflected even in cases
where the other evidence-based components are uninformative.

The final authority score is the weighted sum of the six components, with weights that
reflect their relative diagnostic value: interventions and acts carry the highest weights,
given that they represent direct topical engagement, while profession, education, and role
carry lower weights as background priors. The composite score is percentile-normalised
across the full deputy population, making it comparable across queries.

\paragraph{Stage 5: Coverage Enforcement.}
Before generation proceeds, the system audits the retained evidence set to verify that
each parliamentary group is represented by at least one high-quality fragment. If any
group falls below the coverage threshold, the system selects the best available representative
for that group and ensures it is included in the generation context. The enforcement is
not merely a retrieval correction but a generative constraint: the generation module is
subsequently instructed to produce a dedicated narrative section for each group present
in the evidence, preventing the LLM from implicitly silencing groups whose evidence is
sparse but present.

\paragraph{Stage 6: Ideological Projection.}
The retained evidence fragments are projected into a two-dimensional ideological space
using a weighted variant of principal component analysis (PCA) applied to their dense
embeddings. The first two principal components capture the dominant axes of semantic
variation among the fragments, which---in the context of parliamentary debate---typically
correspond to ideological dimensions interpretable as left-right positioning and a secondary
axis of salience or issue framing.

Each fragment is assigned a two-dimensional coordinate, and these coordinates are
aggregated per parliamentary group using kernel density estimation to produce a smooth
positional estimate. The resulting compass representation is not a fixed ideological map
derived from external political science classifications; it is an \textit{empirical,
query-specific} projection that reflects the actual distribution of argumentative
positions in the retrieved evidence. This approach allows the system to capture
issue-specific realignments---cases in which the standard left-right ordering breaks down
for particular policy domains---that would be invisible to a model based on fixed party
positions.

The stability of the projection is assessed through the ratio of variance explained by the
second principal component relative to the first. When this ratio falls below a threshold,
the system collapses to a one-dimensional representation, signaling to the presentation
layer that the evidence does not support a well-defined second axis.

\paragraph{Stage 7: Multi-Stage Generation.}
Generation is structured as a four-step pipeline of large language model invocations, each
with a distinct functional scope.

In the first step, a \textit{claim analyst} receives the query and the full evidence set
and produces a decomposition of the query into atomic claims---the elementary
propositional questions that a complete answer would need to address. This decomposition
serves as a shared scaffold for the subsequent steps, ensuring that all party-specific
sections address the same epistemic agenda rather than pursuing divergent narrative threads.

In the second step, a \textit{sectional writer} produces, for each parliamentary group, a
prose narrative that addresses the claim set using the evidence attributable to that group.
Each section is generated independently, with the LLM constrained to cite only evidence
fragments from the corresponding group. Citations are embedded as structured placeholders
within the generated text, registering the identifier of the source chunk at the point of
use. The citation registry accumulates all citation expectations across all sections,
forming the authoritative record against which subsequent stages are validated.

In the third step, a \textit{narrative integrator} receives all party-specific sections and
synthesises them into a coherent unified discourse that preserves the internal logic and
attribution of each section. The integrator is explicitly instructed to retain all citation
placeholders; a guard mechanism verifies that the citation count in the integrated output
matches the registry, and a repair procedure injects missing section content as fallback
paragraphs when the integrator inadvertently drops citations during synthesis.

The fourth step---the \textit{citation surgeon}---is not a generative operation but a
deterministic extraction procedure. It traverses the integrated text, resolves each
citation placeholder by retrieving the corresponding chunk's offset pair from the data
layer, and extracts the verbatim quoted string directly from the raw speech text. The
extracted string is formatted according to the system's citation template and inserted in
place of the placeholder. Because extraction is offset-based, this step is guaranteed to
produce a literal substring of the original source, independent of any semantic
interpretation.

\paragraph{Stage 8: Deterministic Citation Insertion.}
The citation surgeon constitutes the final enforcement point of the faithfulness invariant.
Its operation is deliberately non-generative: the surgeon does not ask the LLM to produce
a representative quotation; it computes a quotation by extracting a contiguous character
span from the authoritative source. When the sectional writer has embedded an inline
quotation as part of a citation placeholder---indicating the approximate span of the
intended quotation---the surgeon verifies that this inline string is a literal substring
of the source before accepting it as the citation text. If verification fails, the surgeon
falls back to a semantically-guided sentence extraction that selects the most query-relevant
sentences within the chunk's boundaries, again producing a verbatim extraction rather than
a paraphrase. The final output carries, for each citation, the extracted quote, the
speaker's name, party affiliation, session date, and a unique evidence identifier that
enables independent verification.

\paragraph{Stage 9: Multi-View Structured Answer.}
The final output of the pipeline is a structured document in which the generated narrative
is decomposed into labelled per-party sections, each accompanied by its verified citations.
The document is supplemented by the ideological compass analysis produced in Stage 6, the
authority profiles of the cited speakers, and summary metadata about the coverage of
parliamentary groups in the evidence base. This structured format is designed to support
both direct reading---in which the user reads the synthesised narrative---and analytical
use---in which the user interrogates the underlying evidence and positioning data to form
their own interpretation of the parliamentary debate.
